\chapter{Divisibilità nei numeri interi}\label{cap:divisibilita}

In questa parte studiamo l'aritmetica dei numeri interi con gli strumenti costruiti finora. In questo capitolo: divisori e
multipli, la divisione con resto, il massimo comun divisore e l'algoritmo per calcolarlo, il teorema di Bézout, i numeri
primi e il \textbf{teorema fondamentale dell'aritmetica}, che dice che ogni intero si scompone in modo essenzialmente unico in
fattori primi.

% ──────────────────────────────────────────────────────────────
\section{Divisori, multipli, elementi associati}\label{sec:divisori}

\dfn[divide]{Divisibilità}{
    Sia $s = \mbN$ oppure $s = \mbZ$ e siano $a, b \in s$. Diciamo che \textbf{$a$ divide $b$} (o che $a$ è un \textbf{divisore} di $b$, o che
    $b$ è un \textbf{multiplo} di $a$) e scriviamo $a \mid b$ se $\exists k \in s\,(b = ak)$. Poniamo
    \[ \mathrm{Div}(a) := \{b \in s \mid b \mid a\}, \qquad \mathrm{Mult}(a) := \{b \in s \mid a \mid b\}. \]
    Il \textbf{valore assoluto} di $n \in \mbZ$ è $|n| := \max\{n, -n\}$. I divisori $\pm 1$ e $\pm n$ di un intero $n$ si dicono \textbf{divisori banali} di $n$.
}

\dfn[associati]{Elementi associati}{
    Due interi $a$ e $b$ si dicono \textbf{associati} se $a \mid b$ e $b \mid a$.
}

\begin{concetto}
\begin{Proposizione}{Proprietà della divisibilità in $\mbZ$}{prop-divisibilita}
    \begin{prerequisiti}
        \dip{Definizione}{def:divide}
        \dip{Definizione}{def:associati}
        \dip{Definizione}{def:unita}
    \end{prerequisiti}
    Siano $a, b, c \in \mbZ$.
    \begin{enumerate}
        \item $a \mid a$; se $a \mid b$ e $b \mid c$, allora $a \mid c$; inoltre $1 \mid a$ e $a \mid 0$.
        \item Se $a \mid b$ e $a \mid c$, allora $a \mid bu + cv$ per ogni $u, v \in \mbZ$: un divisore comune divide ogni \textbf{combinazione}.
        \item $a$ e $b$ sono associati se e solo se $b = a$ oppure $b = -a$.
        \item $a$ e $b$ sono associati se e solo se $\mathrm{Div}(a) = \mathrm{Div}(b)$, se e solo se $\mathrm{Mult}(a) = \mathrm{Mult}(b)$.
    \end{enumerate}
\end{Proposizione}
\begin{dimostrazione}
    \begin{caso}{Punti 1 e 2}
        $a = a \cdot 1$; se $b = ah$ e $c = bk$ allora $c = a(hk)$; $a = 1 \cdot a$; $0 = a \cdot 0$. Se $b = ah$ e $c = ak$, allora
        $bu + cv = a(hu + kv)$.
    \end{caso}
    \begin{caso}{Punto 3}
        $(\leftarrow)$ $a = (-a)(-1)$ e $-a = a(-1)$. $(\rightarrow)$ Siano $b = ah$ e $a = bk$. Se $a = 0$, allora $b = 0 = a$. Se $a \neq 0$, da
        $a = a(hk)$ e dalla cancellabilità di $a$ in $(\mbZ, \cdot)$ si ha $hk = 1$: $h$ è invertibile in $\mbZ$, cioè $h = \pm 1$, e $b = \pm a$.
    \end{caso}
    \begin{caso}{Punto 4}
        $(\rightarrow)$ Se $a \mid b$ e $b \mid a$, per la transitività un divisore di $a$ divide $b$ e viceversa; allo stesso modo per i multipli.
        $(\leftarrow)$ Se $\mathrm{Div}(a) = \mathrm{Div}(b)$, da $a \in \mathrm{Div}(a)$ segue $a \mid b$ e da $b \in \mathrm{Div}(b)$ segue $b \mid a$. Se
        $\mathrm{Mult}(a) = \mathrm{Mult}(b)$, da $b \in \mathrm{Mult}(b) = \mathrm{Mult}(a)$ segue $a \mid b$, e simmetricamente $b \mid a$.
    \end{caso}
\end{dimostrazione}
\end{concetto}

\begin{spiegazione}
    In $\mbZ$ la divisibilità non distingue un numero dal suo opposto: $6$ e $-6$ hanno esattamente gli stessi divisori
    ($\pm 1, \pm 2, \pm 3, \pm 6$) e gli stessi multipli. Per questo $\mid$ su $\mbZ$ non è antisimmetrica (lo abbiamo notato
    negli esempi della Sezione~\ref{sec:proprieta-relazioni}) e, quando cercheremo «il» massimo comun divisore, ne troveremo due, uno opposto dell'altro. Il punto 2 è quello
    che useremo più spesso: ad esempio, siccome $3 \mid 12$ e $3 \mid 21$, allora $3$ divide anche $12 \cdot 5 - 21 \cdot 2 = 18$.
\end{spiegazione}

\dfn[mcd]{Massimo comun divisore, minimo comune multiplo}{
    Sia $s = \mbN$ oppure $\mbZ$ e siano $a, b \in s$. Definiamo
    \[ \mathrm{MCD}(a, b) := \{d \in \mathrm{Div}(a) \cap \mathrm{Div}(b) \mid \forall z \in \mathrm{Div}(a) \cap \mathrm{Div}(b)\,(z \mid d)\}, \]
    \[ \mathrm{mcm}(a, b) := \{d \in \mathrm{Mult}(a) \cap \mathrm{Mult}(b) \mid \forall z \in \mathrm{Mult}(a) \cap \mathrm{Mult}(b)\,(d \mid z)\}: \]
    un massimo comun divisore è un divisore comune \textbf{diviso da tutti} i divisori comuni; un minimo comune multiplo è un
    multiplo comune che \textbf{divide tutti} i multipli comuni.
}

\begin{concetto}
\begin{Proposizione}{I massimi comuni divisori sono associati}{mcd-associati}
    \begin{prerequisiti}
        \dip{Definizione}{def:mcd}
        \dip{Proposizione}{prop:prop-divisibilita}
    \end{prerequisiti}
    Siano $a, b \in \mbZ$. Due elementi di $\mathrm{MCD}(a, b)$ sono associati, e un associato di un elemento di $\mathrm{MCD}(a, b)$ sta in
    $\mathrm{MCD}(a, b)$. Quindi, se $d \in \mathrm{MCD}(a, b)$, allora $\mathrm{MCD}(a, b) = \{d, -d\}$. Lo stesso vale per $\mathrm{mcm}(a, b)$.
\end{Proposizione}
\begin{dimostrazione}
    Se $d, d' \in \mathrm{MCD}(a, b)$, sono entrambi divisori comuni; $d'$ divide $d$ perché $d$ è diviso da tutti i divisori comuni, e
    simmetricamente $d \mid d'$. Se invece $d \in \mathrm{MCD}(a, b)$ e $d'$ è associato a $d$, $d'$ ha gli stessi divisori e multipli di $d$
    (Proposizione~\ref{prop:prop-divisibilita}, punto 4): è ancora un divisore comune ed è diviso da tutti i divisori comuni. L'ultima
    affermazione segue dal punto 3 della stessa proposizione. Per $\mathrm{mcm}$ il ragionamento è identico.
\end{dimostrazione}
\end{concetto}

\begin{spiegazione}
    «Massimo» va inteso rispetto alla divisibilità, non rispetto a $\leq$: si guarda il diagramma di Hasse di $(\mbN, \mid)$.
    Per $a = 12$ e $b = 18$, in $\mbN$ i divisori comuni sono $1, 2, 3, 6$ e tutti dividono $6$: quindi $\mathrm{MCD}_{\mbN}(12, 18) = \{6\}$, e in
    $\mbZ$ $\mathrm{MCD}(12, 18) = \{6, -6\}$.
    \begin{center}
    \begin{tikzpicture}[hasse]
        \node[hv evid2] (1) at (0,0) {$1$}; \node[hv evid2] (2) at (-0.9,1.0) {$2$}; \node[hv evid2] (3) at (0.9,1.0) {$3$};
        \node[hv] (4) at (-1.8,2.0) {$4$}; \node[hv evid2] (6) at (0,2.0) {$6$}; \node[hv] (9) at (1.8,2.0) {$9$};
        \node[hv evid] (12) at (-0.9,3.0) {$12$}; \node[hv evid] (18) at (0.9,3.0) {$18$};
        \draw[hlato] (1) -- (2) -- (4) -- (12); \draw[hlato] (1) -- (3) -- (6) -- (12);
        \draw[hlato] (2) -- (6) -- (18); \draw[hlato] (3) -- (9) -- (18);
        \node[font=\small, align=left, anchor=west] at (2.8,1.5)
            {divisori di $12$ e di $18$ in $(\mbN, \mid)$;\\ in verde i divisori comuni:\\ il «più alto» è $6$, e tutti lo dividono.};
    \end{tikzpicture}
    \end{center}
    Un caso curioso: ogni intero divide $0$, quindi $\mathrm{Div}(0) = \mbZ$ e $\mathrm{MCD}(0, 0) = \{0\}$: in $(\mbN, \mid)$ lo $0$ è il massimo.
\end{spiegazione}

% ──────────────────────────────────────────────────────────────
\section{La divisione euclidea}\label{sec:divisione-euclidea}

\begin{concetto}
\begin{Teorema}{Teorema della divisione euclidea (con resto)}{divisione-euclidea}
    \begin{prerequisiti}
        \dip{Definizione}{def:divide}
        \dip{Teorema}{teo:induzione-seconda}
    \end{prerequisiti}
    Per ogni $n, m \in \mbZ$ con $m \neq 0$ esiste una e una sola coppia $(q, r) \in \mbZ \times \mbN$ tale che
    \[ n = mq + r \qquad\text{e}\qquad 0 \leq r < |m|. \]
    $q$ si chiama \textbf{quoziente} e $r$ \textbf{resto} della divisione di $n$ per $m$.
\end{Teorema}
\begin{dimostrazione}
    \begin{caso}{Esistenza per $n \geq 0$, per induzione di seconda forma su $n$}
        Se $n < |m|$ (in particolare se $n = 0$), basta $q = 0$, $r = n$. Se $n \geq |m|$, allora $0 \leq n - |m| < n$ e per ipotesi induttiva
        $n - |m| = mq_1 + r_1$ con $0 \leq r_1 < |m|$. Quindi $n = mq_1 + |m| + r_1$: se $m > 0$, $n = m(q_1 + 1) + r_1$; se $m < 0$, $|m| = -m$ e
        $n = m(q_1 - 1) + r_1$.
    \end{caso}
    \begin{caso}{Esistenza per $n < 0$}
        Per il caso precedente $-n = mq' + r'$ con $0 \leq r' < |m|$. Se $r' = 0$, $n = m(-q')$. Se $r' > 0$, scriviamo
        $n = -mq' - |m| + (|m| - r')$, con $0 < |m| - r' < |m|$; poi $-mq' - |m| = m(-q' - 1)$ se $m > 0$ e $= m(-q' + 1)$ se $m < 0$.
    \end{caso}
    \begin{caso}{Unicità}
        Se $n = mq_1 + r_1 = mq_2 + r_2$ con $0 \leq r_1, r_2 < |m|$, allora $m(q_1 - q_2) = r_2 - r_1$ e quindi $|m|\,|q_1 - q_2| = |r_2 - r_1| < |m|$.
        Poiché $|m| \neq 0$, si ha $|q_1 - q_2| < 1$, cioè $q_1 = q_2$, e di conseguenza $r_1 = r_2$.
    \end{caso}
\end{dimostrazione}
\end{concetto}

\begin{spiegazione}
    Sulla retta dei numeri i multipli di $m$ formano una «scala» con gradini lunghi $|m|$: ogni intero $n$ sta su un gradino o tra due
    gradini consecutivi. $mq$ è il gradino \textbf{subito a sinistra} di $n$ (o $n$ stesso), e $r$ è quanto manca per arrivare a $n$.
    Per $n = -21$ e $m = 4$:
    \begin{center}
    \begin{tikzpicture}[x=0.55cm]
        \draw[-{Stealth}, black!50] (-27,0) -- (2,0);
        \foreach \k [evaluate=\k as \v using int(4*\k)] in {-6,...,0} { \draw[primary!80!black, line width=1pt] (4*\k,-0.15) -- (4*\k,0.15); \node[font=\scriptsize, primary!80!black] at (4*\k,0.45) {$\v$}; }
        \node[circle, fill=orange!80!black, inner sep=1.5pt] at (-21,0) {};
        \node[font=\scriptsize, orange!80!black] at (-21,-0.45) {$n = -21$};
        \draw[decorate, decoration={brace, amplitude=3pt, mirror}] (-24,-0.75) -- node[below=3pt, font=\scriptsize] {$r = 3$} (-21,-0.75);
        \node[font=\scriptsize, align=left, anchor=west] at (-15,-0.8) {$-21 = 4 \cdot (-6) + 3$: \ $q = -6$, $r = 3$\\ (il resto è sempre $\geq 0$!)};
    \end{tikzpicture}
    \end{center}
    Attenzione al segno: $-21 = 4 \cdot (-5) - 1$ non va bene, perché il resto deve essere tra $0$ e $|m| - 1$. Altri esempi:
    $21 = 4 \cdot 5 + 1$, \ $21 = (-4)(-5) + 1$, \ $-21 = (-4) \cdot 6 + 3$, \ $11 = 2 \cdot 5 + 1$, \ $10 = 5 \cdot 2 + 0$.
\end{spiegazione}

% ──────────────────────────────────────────────────────────────
\section{L'algoritmo delle divisioni successive}\label{sec:euclide}

Poiché $\mathrm{MCD}(a, b) = \{d, -d\}$, per trovarlo basta trovare $d \geq 0$; e poiché $a$ e $-a$ hanno gli stessi divisori, possiamo
supporre $a, b \in \mbN$. L'idea per calcolarlo è che dividere con resto \textbf{non cambia i divisori comuni}.

\begin{concetto}
\begin{Lemma}{La divisione non cambia i divisori comuni}{mcd-resto}
    \begin{prerequisiti}
        \dip{Proposizione}{prop:prop-divisibilita}
        \dip{Definizione}{def:mcd}
    \end{prerequisiti}
    Se $a = bq + r$ con $a, b, q, r \in \mbZ$, allora $\mathrm{Div}(a) \cap \mathrm{Div}(b) = \mathrm{Div}(b) \cap \mathrm{Div}(r)$, e quindi
    $\mathrm{MCD}(a, b) = \mathrm{MCD}(b, r)$. Inoltre $\mathrm{MCD}(a, 0) = \{a, -a\}$.
\end{Lemma}
\begin{dimostrazione}
    Se $z \mid a$ e $z \mid b$, allora $z \mid a - bq = r$; se $z \mid b$ e $z \mid r$, allora $z \mid bq + r = a$ (Proposizione~\ref{prop:prop-divisibilita},
    punto 2). Avendo gli stessi divisori comuni, le due coppie hanno gli stessi massimi comuni divisori. Infine, ogni intero divide $0$,
    quindi i divisori comuni di $a$ e $0$ sono i divisori di $a$, e quelli divisi da tutti sono $a$ e $-a$.
\end{dimostrazione}
\end{concetto}

\begin{spiegazione}
    Ad esempio $111 = 17 \cdot 6 + 9$: un numero che divide $111$ e $17$ divide anche $9 = 111 - 6 \cdot 17$, e uno che divide $17$ e $9$
    divide anche $111$. Quindi cercare l'MCD di $111$ e $17$ è come cercare quello di $17$ e $9$: numeri più piccoli.
\end{spiegazione}

\begin{concetto}
\begin{Teorema}{Algoritmo delle divisioni successive (di Euclide)}{euclide}
    \begin{prerequisiti}
        \dip{Teorema}{teo:divisione-euclidea}
        \dip{Lemma}{lem:mcd-resto}
        \dip{Definizione}{def:ben-ordinato}
    \end{prerequisiti}
    Siano $a, b \in \mbN$. Si ripete: se $b = 0$, ci si ferma; altrimenti si divide $a$ per $b$, $a = bq + r$, e si sostituisce la coppia $(a, b)$
    con $(b, r)$. Il procedimento termina, e quando si ferma, con la coppia $(a', 0)$, si ha $a' \in \mathrm{MCD}(a, b)$.
\end{Teorema}
\begin{dimostrazione}
    Posto $r_0 = b$, l'algoritmo produce i resti $r_1 > r_2 > \dots \geq 0$: ogni resto è minore del divisore usato, che è il resto precedente.
    Una successione strettamente decrescente di numeri naturali non può essere infinita (l'insieme dei suoi termini avrebbe un minimo, ma
    il termine successivo sarebbe più piccolo), quindi a un certo passo il resto è $0$. Per il Lemma~\ref{lem:mcd-resto}, ad ogni passo
    $\mathrm{MCD}(a, b)$ resta lo stesso, e alla fine $\mathrm{MCD}(a, b) = \mathrm{MCD}(a', 0) = \{a', -a'\}$.
\end{dimostrazione}
\end{concetto}

\begin{spiegazione}
    Ecco l'algoritmo come diagramma di flusso (per $a, b \in \mbN$ non entrambi nulli) e un esempio passo passo con $a = 111$, $b = 17$:
    \begin{center}
    \begin{tikzpicture}[node distance=0.9cm, every node/.style={font=\small},
        blocco/.style={rectangle, draw=primary!70!black, fill=primary!7, rounded corners=2pt, minimum height=6mm, inner xsep=6pt},
        scelta/.style={diamond, aspect=2.2, draw=orange!80!black, fill=orange!12, inner sep=1pt},
        estremo/.style={ellipse, draw=black!60, fill=black!5, minimum height=6mm},
        flusso/.style={-{Stealth[length=2mm]}, line width=0.7pt, draw=black!70}]
        \node[estremo] (inizio) at (0,0) {Inizio};
        \node[scelta] (a0) at (0,-1.2) {$a = 0$?};
        \node[scelta] (b0) at (0,-2.7) {$b = 0$?};
        \node[blocco] (de) at (0,-4.1) {$(q, r) :=$ divisione di $a$ per $b$};
        \node[blocco] (sc) at (0,-5.2) {$(a, b) := (b, r)$};
        \node[blocco] (ra) at (3.6,-1.2) {$\mathrm{MCD} := \{b\}$};
        \node[blocco] (rb) at (3.6,-2.7) {$\mathrm{MCD} := \{a\}$};
        \node[estremo] (fine) at (5.9,-1.95) {Fine};
        \draw[flusso] (inizio) -- (a0); \draw[flusso] (a0) -- node[right] {no} (b0); \draw[flusso] (b0) -- node[right] {no} (de);
        \draw[flusso] (de) -- (sc); \draw[flusso] (sc.west) -- ++(-1.9,0) |- (b0.west);
        \draw[flusso] (a0) -- node[above] {sì} (ra); \draw[flusso] (b0) -- node[above] {sì} (rb);
        \draw[flusso] (ra) -| (fine); \draw[flusso] (rb) -| (fine);
        \node[anchor=west, align=left] at (7.0,-2.6) {
            \renewcommand{\arraystretch}{1.15}
            \begin{tabular}{c|c|l}
                passo & $(a, b)$ & divisione \\ \hline
                $1$ & $(111, 17)$ & $111 = 17 \cdot 6 + 9$ \\
                $2$ & $(17, 9)$ & $17 = 9 \cdot 1 + 8$ \\
                $3$ & $(9, 8)$ & $9 = 8 \cdot 1 + 1$ \\
                $4$ & $(8, 1)$ & $8 = 1 \cdot 8 + 0$ \\
                fine & $(1, 0)$ & $\mathrm{MCD} = \{1\}$ in $\mbN$ \\
            \end{tabular}};
    \end{tikzpicture}
    \end{center}
    I resti $9, 8, 1, 0$ scendono sempre, quindi l'algoritmo non può andare avanti all'infinito. L'ultimo resto non nullo, $1$, è il
    massimo comun divisore: $\mathrm{MCD}(111, 17) = \{1, -1\}$ in $\mbZ$.
\end{spiegazione}

% ──────────────────────────────────────────────────────────────
\section{Il teorema di Bézout}\label{sec:bezout}

\begin{concetto}
\begin{Teorema}{Teorema di Bézout}{bezout}
    \begin{prerequisiti}
        \dip{Teorema}{teo:euclide}
        \dip{Proposizione}{prop:mcd-associati}
        \dip{Teorema}{teo:induzione-seconda}
    \end{prerequisiti}
    Siano $a, b \in \mbZ$ non entrambi nulli e sia $d \in \mathrm{MCD}(a, b)$. Allora esistono $u, v \in \mbZ$ tali che $d = au + bv$.
\end{Teorema}
\begin{dimostrazione}
    Se vale per $d$ vale per $-d = a(-u) + b(-v)$; e poiché $a = \pm |a|$, $b = \pm |b|$, una combinazione di $|a|$ e $|b|$ è anche una
    combinazione di $a$ e $b$. Possiamo quindi supporre $a, b \in \mbN$ e $d \geq 0$. Se $b = 0$, $d = a = a \cdot 1 + b \cdot 0$. Altrimenti
    consideriamo i resti dell'algoritmo di Euclide, $r_{-1} = a$, $r_0 = b$, $r_1, \dots, r_{t - 1}$, $r_t = 0$, con $r_{i - 2} = r_{i - 1}q_i + r_i$ e
    $d = r_{t - 1}$. Dimostriamo di più: \textbf{ogni resto è una combinazione di $a$ e $b$}. Per assurdo sia $k$ il minimo indice per cui
    $r_k$ non lo è. Non è $k = -1$ né $k = 0$ ($a = a \cdot 1 + b \cdot 0$, $b = a \cdot 0 + b \cdot 1$). Allora $r_{k - 2} = au' + bv'$ e
    $r_{k - 1} = au'' + bv''$, e
    \[ r_k = r_{k - 2} - r_{k - 1}q_k = a(u' - u''q_k) + b(v' - v''q_k), \]
    che è una combinazione di $a$ e $b$: assurdo. In particolare $d = r_{t - 1}$ è una combinazione di $a$ e $b$.
\end{dimostrazione}
\end{concetto}

\begin{spiegazione}
    La dimostrazione è anche un \textbf{algoritmo} (algoritmo delle divisioni successive \textbf{esteso}): si scrive accanto a ogni resto
    come combinazione di $a$ e $b$, usando le due righe precedenti. Per $a = 111$, $b = 17$:
    \begin{center}
    \renewcommand{\arraystretch}{1.25}
    \begin{tabular}{l|c|cc|l}
        \rowcolor{primary!12} riga & resto & $u$ & $v$ & come si ottiene \\
        $r_{-1}$ & $111$ & $1$ & $0$ & $a$ \\
        $r_0$ & $17$ & $0$ & $1$ & $b$ \\
        $r_1$ & $9$ & $1$ & $-6$ & $r_{-1} - 6\,r_0$ \quad ($111 = 17 \cdot 6 + 9$) \\
        $r_2$ & $8$ & $-1$ & $7$ & $r_0 - 1\,r_1$ \quad ($17 = 9 \cdot 1 + 8$) \\
        $r_3$ & $1$ & $2$ & $-13$ & $r_1 - 1\,r_2$ \quad ($9 = 8 \cdot 1 + 1$) \\
    \end{tabular}
    \end{center}
    Ogni riga si ottiene sottraendo dalla penultima l'ultima moltiplicata per il quoziente, \textbf{anche nelle colonne $u$ e $v$}.
    Risultato: $1 = 111 \cdot 2 + 17 \cdot (-13)$. Verifica: $222 - 221 = 1$.
\end{spiegazione}

% ──────────────────────────────────────────────────────────────
\section{Numeri coprimi, lemma di Euclide, numeri primi}\label{sec:primi}

\dfn[coprimi]{Numeri coprimi}{
    Due interi $a$ e $b$ si dicono \textbf{coprimi} se $1 \in \mathrm{MCD}(a, b)$, cioè se i loro unici divisori comuni sono $1$ e $-1$.
}

\begin{concetto}
\begin{Teorema}{Lemma di Euclide}{lemma-euclide}
    \begin{prerequisiti}
        \dip{Definizione}{def:coprimi}
        \dip{Teorema}{teo:bezout}
        \dip{Proposizione}{prop:prop-divisibilita}
    \end{prerequisiti}
    Siano $a, b, c \in \mbZ$. Se $a$ e $b$ sono coprimi e $a \mid bc$, allora $a \mid c$.
\end{Teorema}
\begin{dimostrazione}
    Per ipotesi $bc = ah$ per qualche $h$. Per il teorema di Bézout esistono $u, v$ con $1 = au + bv$. Moltiplicando per $c$:
    \[ c = acu + bcv = acu + ahv = a(cu + hv), \]
    quindi $a \mid c$.
\end{dimostrazione}
\end{concetto}

\begin{spiegazione}
    Se $a$ divide un prodotto ma non ha niente in comune con il primo fattore, deve «stare tutto» nel secondo. Ad esempio $4 \mid 3 \cdot 8$
    e $\mathrm{MCD}(4, 3) = \{\pm 1\}$, quindi $4 \mid 8$. Senza l'ipotesi di coprimalità è falso: $6 \mid 4 \cdot 9 = 36$, ma $6 \nmid 4$ e $6 \nmid 9$
    ($6$ si «divide» tra i due fattori: il $2$ va nel $4$, il $3$ nel $9$).
\end{spiegazione}

\dfn[primo]{Numero primo, elemento irriducibile}{
    Un intero $p \notin \{-1, 0, 1\}$ si dice
    \begin{itemize}
        \item \textbf{primo} se $\forall a, b \in \mbZ\,(p \mid ab \rightarrow p \mid a \vee p \mid b)$;
        \item \textbf{irriducibile} se ha solo i divisori banali $\pm 1$, $\pm p$.
    \end{itemize}
}

\begin{concetto}
\begin{Teorema}{In $\mbZ$ primi e irriducibili coincidono}{primi-irriducibili}
    \begin{prerequisiti}
        \dip{Definizione}{def:primo}
        \dip{Teorema}{teo:lemma-euclide}
        \dip{Proposizione}{prop:prop-divisibilita}
    \end{prerequisiti}
    Sia $p \in \mbZ \setminus \{-1, 0, 1\}$. Allora $p$ è primo se e solo se è irriducibile.
\end{Teorema}
\begin{dimostrazione}
    \begin{caso}{$(\rightarrow)$ Primo $\Rightarrow$ irriducibile}
        Sia $a$ un divisore di $p$: $p = ab$. Allora $p \mid ab$ e, essendo primo, $p \mid a$ oppure $p \mid b$. Se $p \mid a$, $a = pk$ e
        $p = pkb$, da cui $kb = 1$ ($p \neq 0$ è cancellabile): $b = \pm 1$ e $a = \pm p$. Se $p \mid b$, allo stesso modo $a = \pm 1$.
        Quindi i divisori di $p$ sono solo quelli banali.
    \end{caso}
    \begin{caso}{$(\leftarrow)$ Irriducibile $\Rightarrow$ primo}
        Sia $p \mid ab$ e supponiamo $p \nmid a$. I divisori comuni di $p$ e $a$ sono divisori di $p$, quindi tra $\pm 1, \pm p$; ma $\pm p$ non
        dividono $a$. Quindi $p$ e $a$ sono coprimi e, per il lemma di Euclide, $p \mid b$.
    \end{caso}
\end{dimostrazione}
\end{concetto}

\begin{spiegazione}
    «Irriducibile» dice che $p$ non si spezza in fattori più piccoli; «primo» dice che quando divide un prodotto divide uno dei fattori.
    In $\mbZ$ sono la stessa cosa. Ad esempio $7$ è irriducibile, e infatti se $7 \mid ab$ allora $7$ divide $a$ o $b$; $6$ non è irriducibile
    ($6 = 2 \cdot 3$) e infatti non è primo: $6 \mid 2 \cdot 3$ ma $6 \nmid 2$ e $6 \nmid 3$. I primi positivi sono $2, 3, 5, 7, 11, 13, \dots$; anche $-2, -3, \dots$
    sono primi.
\end{spiegazione}

% ──────────────────────────────────────────────────────────────
\section{Il teorema fondamentale dell'aritmetica}\label{sec:tfa}

\begin{concetto}
\begin{Lemma}{Un primo che divide un prodotto divide un fattore}{primo-prodotto}
    \begin{prerequisiti}
        \dip{Definizione}{def:primo}
        \dip{Corollario}{cor:induzione-uso}
    \end{prerequisiti}
    Se $p$ è primo e $p \mid a_1 a_2 \cdots a_s$ con $s \geq 1$, allora $p \mid a_j$ per qualche $j \in \{1, \dots, s\}$.
\end{Lemma}
\begin{dimostrazione}
    Per induzione su $s$. Per $s = 1$ è ovvio. Se $p \mid (a_1 \cdots a_s) a_{s + 1}$, per definizione di primo $p \mid a_1 \cdots a_s$ (e si conclude con
    l'ipotesi induttiva) oppure $p \mid a_{s + 1}$.
\end{dimostrazione}
\end{concetto}

\begin{concetto}
\begin{Teorema}{Teorema fondamentale dell'aritmetica}{tfa}
    \begin{prerequisiti}
        \dip{Teorema}{teo:primi-irriducibili}
        \dip{Lemma}{lem:primo-prodotto}
        \dip{Teorema}{teo:induzione-seconda}
        \dip{Corollario}{cor:induzione-uso}
    \end{prerequisiti}
    Sia $m \in \mbZ \setminus \{-1, 0, 1\}$. Allora:
    \begin{enumerate}
        \item (\textbf{esistenza}) esistono primi $p_1, \dots, p_r$ tali che $m = p_1 p_2 \cdots p_r$;
        \item (\textbf{essenziale unicità}) se anche $m = q_1 \cdots q_s$ con i $q_j$ primi, allora $r = s$ ed esiste una funzione biettiva
              $\sigma \colon I_r \to I_r$ tale che $p_i = \pm q_{\sigma(i)}$ per ogni $i \in I_r$.
    \end{enumerate}
\end{Teorema}
\begin{dimostrazione}
    \begin{caso}{Esistenza per $m \geq 2$, per induzione di seconda forma}
        $2$ è primo, perché i suoi unici divisori positivi sono $1$ e $2$. Sia $m > 2$ e supponiamo la tesi per ogni $k$ con $2 \leq k < m$.
        Se $m$ è primo abbiamo finito ($r = 1$). Altrimenti $m$ non è irriducibile e ha un divisore non banale: $m = ab$ con $a, b \in \mbN$ e
        $1 < a, b < m$. Per ipotesi induttiva $a$ e $b$ sono prodotti di primi, quindi lo è anche $m = ab$.
    \end{caso}
    \begin{caso}{Esistenza per $m \leq -2$}
        $-m = p_1 \cdots p_r$, quindi $m = (-p_1) p_2 \cdots p_r$, e $-p_1$ è primo (ha gli stessi divisori di $p_1$).
    \end{caso}
    \begin{caso}{Essenziale unicità, per induzione su $r$}
        Sia $p_1 \cdots p_r = q_1 \cdots q_s$. Poiché $p_r$ divide il secondo membro, per il Lemma~\ref{lem:primo-prodotto} divide qualche $q_j$;
        rinumerando i $q$ possiamo supporre $p_r \mid q_s$. Poiché $q_s$ è irriducibile e $p_r \neq \pm 1$, si ha $q_s = \pm p_r$. Cancellando
        $p_r$ ($\neq 0$) si ottiene $p_1 \cdots p_{r - 1} = \pm q_1 \cdots q_{s - 1}$. Se $r = 1$, il primo membro è $1$, quindi anche $s = 1$ (nessun primo
        divide $\pm 1$). Se $r > 1$, allo stesso modo $s > 1$ (altrimenti $p_1$ dividerebbe $\pm 1$), e si applica l'ipotesi induttiva (assorbendo il segno in $q_1$): $r - 1 = s - 1$ e c'è una biezione
        $\tau$ di $I_{r - 1}$ con $p_i = \pm q_{\tau(i)}$; estendendo $\tau$ con $\sigma(r) = r$ (dopo la rinumerazione) si ottiene $\sigma$.
    \end{caso}
\end{dimostrazione}
\end{concetto}

\begin{spiegazione}
    \textbf{Esistenza}: si spezza il numero finché tutti i pezzi sono primi; l'induzione garantisce che lo spezzare finisca, perché i
    pezzi sono sempre più piccoli. Ecco l'«albero» di $360$:
    \begin{center}
    \begin{tikzpicture}[level distance=0.85cm, every node/.style={font=\small},
                        level 1/.style={sibling distance=2.6cm}, level 2/.style={sibling distance=1.3cm}, level 3/.style={sibling distance=0.8cm},
                        edge from parent/.style={draw=black!50}]
        \node {$360$}
            child {node {$12$} child {node {$4$} child {node[hv evid] {$2$}} child {node[hv evid] {$2$}}} child {node[hv evid] {$3$}}}
            child {node {$30$} child {node[hv evid] {$2$}} child {node {$15$} child {node[hv evid] {$3$}} child {node[hv evid] {$5$}}}};
        \node[anchor=west, font=\small] at (3.2,-1.4) {$360 = 2 \cdot 2 \cdot 3 \cdot 2 \cdot 3 \cdot 5 = 2^3 \cdot 3^2 \cdot 5$};
    \end{tikzpicture}
    \end{center}
    \textbf{Unicità}: spezzando in modo diverso ($360 = 10 \cdot 36$, ad esempio) si ottengono gli stessi primi, al più in un altro ordine e
    con segni diversi, come in $12 = 2 \cdot 2 \cdot 3 = 3 \cdot (-2) \cdot (-2)$. La dimostrazione «stacca» un primo alla volta: $p_r$ deve comparire
    anche nell'altra scomposizione (è qui che serve che sia \textbf{primo}), lo si cancella da entrambe e si continua.
\end{spiegazione}

\begin{concetto}
\begin{Corollario}{Esistono infiniti numeri primi (Euclide)}{infiniti-primi}
    \begin{prerequisiti}
        \dip{Teorema}{teo:tfa}
        \dip{Proposizione}{prop:prop-divisibilita}
    \end{prerequisiti}
    L'insieme dei numeri primi non è finito.
\end{Corollario}
\begin{dimostrazione}
    Per assurdo siano $p_1, \dots, p_t$ tutti i primi positivi e sia $m = p_1 \cdots p_t + 1 \geq 2$. Per il Teorema~\ref{teo:tfa} qualche primo
    positivo $p$ divide $m$, e $p = p_i$ per qualche $i$. Ma $p_i$ divide anche $p_1 \cdots p_t$, quindi divide la differenza
    $m - p_1 \cdots p_t = 1$: assurdo.
\end{dimostrazione}
\end{concetto}

\begin{spiegazione}
    Data una lista finita di primi, il numero $m = (\text{prodotto della lista}) + 1$ lascia resto $1$ nella divisione per ciascuno di essi:
    quindi i suoi fattori primi non sono nella lista. Con la lista $2, 3, 5$: $m = 31$, che è un primo nuovo. (Attenzione: $m$ non è sempre
    primo: con $2, 3, 5, 7, 11, 13$ si ha $m = 30031 = 59 \cdot 509$; ma i suoi fattori primi sono comunque nuovi.)
\end{spiegazione}

Dalla scomposizione in fattori primi si leggono tutti i divisori di un numero, e quindi anche massimo comun divisore e minimo
comune multiplo.

\begin{concetto}
\begin{Proposizione}{Divisori, MCD e mcm dalla scomposizione}{divisori-scomposizione}
    \begin{prerequisiti}
        \dip{Teorema}{teo:tfa}
        \dip{Lemma}{lem:somma-prodotto}
        \dip{Definizione}{def:mcd}
    \end{prerequisiti}
    Sia $n = p_1^{e_1} \cdots p_r^{e_r}$ con $p_1, \dots, p_r$ primi positivi distinti ed $e_i \geq 1$. Allora:
    \begin{enumerate}
        \item i divisori positivi di $n$ sono esattamente i numeri $p_1^{f_1} \cdots p_r^{f_r}$ con $0 \leq f_i \leq e_i$; sono $(e_1 + 1) \cdots (e_r + 1)$, e
              i divisori in $\mbZ$ sono il doppio;
        \item se $a = p_1^{e_1} \cdots p_r^{e_r}$ e $b = p_1^{f_1} \cdots p_r^{f_r}$ (con esponenti $\geq 0$), allora
              $d = \prod p_i^{\min(e_i, f_i)} \in \mathrm{MCD}(a, b)$, $\ell = \prod p_i^{\max(e_i, f_i)} \in \mathrm{mcm}(a, b)$ e $d\ell = ab$.
    \end{enumerate}
\end{Proposizione}
\begin{dimostrazione}
    \begin{caso}{Punto 1}
        Ogni $p_1^{f_1} \cdots p_r^{f_r}$ con $f_i \leq e_i$ divide $n$ (il quoziente è $\prod p_i^{e_i - f_i}$). Viceversa, se $c > 0$ e $n = ck$, scomponendo
        $c$ e $k$ in primi positivi si ottiene una scomposizione di $n$; per l'unicità i primi di $c$ sono tra i $p_i$, ognuno con esponente
        al più $e_i$. Per l'unicità, esponenti diversi danno numeri diversi; quindi i divisori positivi corrispondono alle $r$-uple
        $(f_1, \dots, f_r)$, che sono $(e_1 + 1) \cdots (e_r + 1)$ (Lemma~\ref{lem:somma-prodotto}). In $\mbZ$ si aggiungono gli opposti.
    \end{caso}
    \begin{caso}{Punto 2}
        Per il punto 1, i divisori comuni positivi di $a$ e $b$ sono i $\prod p_i^{g_i}$ con $g_i \leq e_i$ e $g_i \leq f_i$, cioè $g_i \leq \min(e_i, f_i)$:
        tutti dividono $d$, che è uno di loro. Allo stesso modo i multipli comuni hanno esponenti $\geq \max(e_i, f_i)$ e sono multipli di $\ell$.
        Infine $\min(e_i, f_i) + \max(e_i, f_i) = e_i + f_i$, quindi $d\ell = ab$.
    \end{caso}
\end{dimostrazione}
\end{concetto}

\begin{spiegazione}
    Un divisore di $12 = 2^2 \cdot 3$ si costruisce scegliendo quante volte usare il $2$ ($0$, $1$ o $2$) e quante il $3$ ($0$ o $1$): $3 \cdot 2 = 6$
    divisori, disposti in una griglia che è proprio il diagramma di Hasse dei divisori di $12$:
    \begin{center}
    \begin{tikzpicture}[hasse]
        \foreach \f/\g/\v in {0/0/1, 1/0/2, 2/0/4, 0/1/3, 1/1/6, 2/1/12} \node[hv] (d\f\g) at (1.5*\f-0.8*\g, 0.9*\f+1.0*\g) {$\v$};
        \draw[hlato] (d00) -- (d10) -- (d20); \draw[hlato] (d01) -- (d11) -- (d21);
        \draw[hlato] (d00) -- (d01); \draw[hlato] (d10) -- (d11); \draw[hlato] (d20) -- (d21);
        \node[font=\small, align=left, anchor=west] at (3.8,1.3)
            {verso destra: un $2$ in più;\\ verso l'alto a sinistra: un $3$ in più.\\[1mm]
             $60 = 2^2 \cdot 3 \cdot 5$ ha $3 \cdot 2 \cdot 2 = 12$ divisori positivi;\\
             $\mathrm{MCD}(12, 18) \ni 2^{\min(2,1)} \cdot 3^{\min(1,2)} = 6$,\\ $\mathrm{mcm}(12, 18) \ni 2^2 \cdot 3^2 = 36$, \ $6 \cdot 36 = 12 \cdot 18$.};
    \end{tikzpicture}
    \end{center}
\end{spiegazione}

% ──────────────────────────────────────────────────────────────
\section{I sottogruppi di $(\mbZ, +)$}\label{sec:sottogruppi-z}

Per ogni $m \in \mbZ$, l'insieme dei multipli $m\mbZ = \langle m \rangle$ è un sottogruppo di $(\mbZ, +)$ (Proposizione~\ref{prop:cicliche}). Non
ce ne sono altri.

\begin{concetto}
\begin{Proposizione}{I sottogruppi di $\mbZ$}{sottogruppi-z}
    \begin{prerequisiti}
        \dip{Definizione}{def:sottostruttura}
        \dip{Teorema}{teo:divisione-euclidea}
        \dip{Definizione}{def:ben-ordinato}
    \end{prerequisiti}
    I sottogruppi di $(\mbZ, +)$ sono tutti e soli gli insiemi $m\mbZ$ con $m \in \mbN$.
\end{Proposizione}
\begin{dimostrazione}
    \facoltativa
    Sia $h$ un sottogruppo. Se $h = \{0\}$, $h = 0\mbZ$. Altrimenti $h$ contiene un elemento non nullo e il suo opposto, quindi un elemento
    positivo; sia $m$ il minimo degli elementi positivi di $h$. I multipli di $m$ stanno in $h$ (somme di $m$ e di $-m$), quindi $m\mbZ \subseteq h$.
    Se $x \in h$, dividiamo: $x = mq + r$ con $0 \leq r < m$; allora $r = x - mq \in h$ e, per la minimalità di $m$, $r = 0$: $x \in m\mbZ$.
\end{dimostrazione}
\end{concetto}

\begin{spiegazione}
    Un sottogruppo di $\mbZ$ è chiuso per somme e differenze. Se contiene il suo più piccolo elemento positivo $m$, contiene tutti i
    multipli di $m$; e se contenesse un numero che non è multiplo di $m$, sottraendo multipli di $m$ si otterrebbe un resto positivo più
    piccolo di $m$, impossibile. Ad esempio il sottogruppo generato da $\{4, 6\}$ contiene $6 - 4 = 2$, quindi è $2\mbZ$.
\end{spiegazione}
